% Options for packages loaded elsewhere
\PassOptionsToPackage{unicode}{hyperref}
\PassOptionsToPackage{hyphens}{url}
%
\documentclass[
]{article}
\usepackage{amsmath,amssymb}
\usepackage{iftex}
\ifPDFTeX
  \usepackage[T1]{fontenc}
  \usepackage[utf8]{inputenc}
  \usepackage{textcomp} % provide euro and other symbols
\else % if luatex or xetex
  \usepackage{unicode-math} % this also loads fontspec
  \defaultfontfeatures{Scale=MatchLowercase}
  \defaultfontfeatures[\rmfamily]{Ligatures=TeX,Scale=1}
\fi
\usepackage{lmodern}
\ifPDFTeX\else
  % xetex/luatex font selection
\fi
% Use upquote if available, for straight quotes in verbatim environments
\IfFileExists{upquote.sty}{\usepackage{upquote}}{}
\IfFileExists{microtype.sty}{% use microtype if available
  \usepackage[]{microtype}
  \UseMicrotypeSet[protrusion]{basicmath} % disable protrusion for tt fonts
}{}
\makeatletter
\@ifundefined{KOMAClassName}{% if non-KOMA class
  \IfFileExists{parskip.sty}{%
    \usepackage{parskip}
  }{% else
    \setlength{\parindent}{0pt}
    \setlength{\parskip}{6pt plus 2pt minus 1pt}}
}{% if KOMA class
  \KOMAoptions{parskip=half}}
\makeatother
\usepackage{xcolor}
\setlength{\emergencystretch}{3em} % prevent overfull lines
\providecommand{\tightlist}{%
  \setlength{\itemsep}{0pt}\setlength{\parskip}{0pt}}
\setcounter{secnumdepth}{-\maxdimen} % remove section numbering
\ifLuaTeX
  \usepackage{selnolig}  % disable illegal ligatures
\fi
\usepackage{bookmark}
\IfFileExists{xurl.sty}{\usepackage{xurl}}{} % add URL line breaks if available
\urlstyle{same}
\hypersetup{
  pdftitle={Research Notebook: Multi-Agent Learning and Liar\textquotesingle s Dice},
  hidelinks,
  pdfcreator={LaTeX via pandoc}}

\title{Research Notebook: Multi-Agent Learning and
Liar\textquotesingle s Dice}
\author{}
\date{}

\begin{document}
\maketitle

\section{Research Notebook: Multi-Agent Learning and Liar's
Dice}\label{research-notebook-multi-agent-learning-and-liars-dice}

\subsection{Public Edition}\label{public-edition}

\textbf{Author:} Mauricio Garcia Villanueva\\
\textbf{Project:} Liar's Dice CFR Lab\\
\textbf{Live demo:} emgarviii.github.io/liars\_dice/

\subsection{Notebook Purpose}\label{notebook-purpose}

This notebook is a cleaned public version of the planning notes behind
the project. The original notebook was exploratory. It had definitions,
copied reminders, rough implementation thoughts, and places where my
understanding changed over time. This edition keeps the research
trajectory but removes clutter and makes the reasoning easier to follow.

The main arc is:

\begin{enumerate}
\def\labelenumi{\arabic{enumi}.}
\tightlist
\item
  Start with imperfect-information games and multi-agent learning.
\item
  Study poker AI and CFR algorithms.
\item
  Try to apply those ideas to Liar's Dice.
\item
  Realize that full classic Liar's Dice is too large for a clean first
  solver.
\item
  Narrow the project into a one-claim challenge abstraction.
\item
  Build, evaluate, and publish the result.
\end{enumerate}

\subsection{Initial Motivation}\label{initial-motivation}

I wanted a research project that connected game theory, machine
learning, and strategic decision making. Poker AI was the obvious
reference point, but I did not want to only summarize poker systems. I
wanted to build something with hidden information that I could explain
visually and eventually publish.

Liar's Dice stood out because it is easy to understand and still
strategically rich. Each player has private dice, and the public claim
creates a decision under uncertainty. The game naturally raises the
questions I was interested in:

\begin{itemize}
\tightlist
\item
  How does an AI reason when it cannot see the full state?
\item
  How does self-play help a strategy improve?
\item
  What does equilibrium mean in a game with hidden information?
\item
  How do we evaluate whether the policy is actually good?
\end{itemize}

\subsection{Concepts I Needed to
Learn}\label{concepts-i-needed-to-learn}

The first concept was the information set. In a perfect-information
game, the player can see the full board. In an imperfect-information
game, the player sees only part of the state. The policy must therefore
be based on what the player knows, not on the true full state.

The second concept was Nash equilibrium. A Nash equilibrium is a stable
set of strategies where no player can improve by changing strategy
alone. I originally treated equilibrium as a stronger and more automatic
claim than I should have. The public version of the project is more
careful. It talks about sampled CFR+ training and best-response
diagnostics, but it does not claim that the published policy is proven
to be at equilibrium.

The third concept was CFR. Counterfactual Regret Minimization repeatedly
updates action probabilities by asking which actions would have
performed better at an information set. CFR+ modifies that process by
clipping negative regret, which can improve convergence behavior in
practice.

The fourth concept was exploitability. A policy can beat simple
opponents but still be weak against a targeted opponent. This became one
of the most important lessons of the project.

\subsection{Early Direction}\label{early-direction}

The early project vision was broader than the final public demo. I
wanted to build something close to classic Liar's Dice, where players
raise claims until someone calls. That version is more familiar and more
realistic.

The problem was the game tree. Once repeated raises and histories are
included, the number of information sets grows quickly. A clean solver
would need more careful abstraction, stronger training, and more
evaluation. That was possible as future work, but it was too much for
the first publishable version.

This is where the project changed from ``solve Liar's Dice'' to ``build
a tractable imperfect-information lab inspired by Liar's Dice.'' That
was a better engineering decision because it created a version I could
actually validate and explain.

\subsection{The Simplified Game}\label{the-simplified-game}

The final public research mode uses one claim and one response.

The claimant says something like ``two of face six.'' The claim refers
to both players' dice combined. The responder sees only their own dice
and must choose Believe or Challenge.

This keeps the hidden-information structure:

\begin{itemize}
\tightlist
\item
  the claimant does not know the responder's dice,
\item
  the responder does not know the claimant's dice,
\item
  the claim may be true, false, or uncertain from one player's view.
\end{itemize}

It removes the classic raise loop. That is a real limitation, but it
also makes the project explainable enough for a public portfolio demo.

\subsection{Implementation Notes}\label{implementation-notes}

The implementation split the project into two layers.

The Python layer is the research pipeline. It defines the simplified
rules, generates game states, trains policies, validates distributions,
simulates matches, and exports JSON artifacts.

The TypeScript layer is the public interface. It loads the exported
policy and metrics files, renders the playable game, explains the
decision guide, and exposes the method, results, classic comparison, and
paper archive pages.

This split ended up being one of the best architectural decisions. It
means GitHub Pages can host the project without a backend. The AI policy
is frozen and reproducible rather than retrained in the browser.

\subsection{Policy Key Lesson}\label{policy-key-lesson}

An early weakness was that the policy did not condition strongly enough
on changing dice counts. A strategy that makes sense at five dice each
may not make sense later in the match when one side has fewer dice.

The updated public policy is remaining-dice-aware. Its keys include
public dice counts and private hands. That improved benchmark
performance because the policy could adapt to match state instead of
reusing a fixed starting-state strategy.

This was an important engineering lesson: the representation of state is
not just a detail. It directly controls what the policy can learn.

\subsection{Evaluation Notes}\label{evaluation-notes}

At first, it was tempting to judge the AI by whether it felt smart in
the browser. That was not enough. A playable demo is useful for
communication, but evaluation needs repeatable metrics.

The benchmark suite compares the policy against several opponent
profiles:

\begin{itemize}
\tightlist
\item
  random behavior,
\item
  skeptical responses,
\item
  threshold-based responses,
\item
  truth-biased claims.
\end{itemize}

The updated policy performs much better across these profiles than the
baseline. The benchmark average rose from 45.3 percent to 73.0 percent.

The best-response-style diagnostic made the story more honest. It showed
that the policy still has exploitable pressure, especially when a
stronger responder targets the claimant policy. That result changed how
the project should be presented. The correct claim is not ``the AI is
optimal.'' The correct claim is ``the AI improved against benchmarks,
and the evaluation exposed where it remains weak.''

\subsection{Website Notes}\label{website-notes}

The original project was not easy enough to show publicly. It had papers
and code, but not a clean path for someone to understand the work
quickly.

The public website solves that by giving different readers different
depths:

\begin{itemize}
\tightlist
\item
  The homepage explains the project and lets people play.
\item
  The method page explains the research engineering pipeline.
\item
  The results page explains the metrics and limitations.
\item
  The classic page shows the familiar Liar's Dice loop for comparison.
\item
  The papers page preserves the original course documents and adds
  polished public editions.
\end{itemize}

This matters because portfolio projects need more than implementation.
They need a clear story.

\subsection{What I Would Improve Next}\label{what-i-would-improve-next}

The biggest technical improvement would be to train directly against
stronger best-response pressure. That would move the project from
benchmark improvement toward robustness.

The second improvement would be to extend the solver toward the classic
raise-and-challenge game. That would require a larger state
representation and more careful abstraction.

The third improvement would be better convergence reporting. The current
checkpoints are useful, but a stronger research version would show more
formal exploitability trends across training.

The fourth improvement would be richer visualization. The site could
show policy behavior by information set after a round is revealed, while
still avoiding hidden-hand leakage during play.

\subsection{Final Reflection}\label{final-reflection}

The most useful part of the project was learning to separate ambition
from evidence. It is easy to say ``CFR finds equilibrium'' in a general
discussion. It is harder, and more valuable, to say exactly what my code
trained, exactly what my tests measured, and exactly where the result
still falls short.

That is why the public artifact is stronger than the original course
submission. It does not just preserve the work. It makes the work
inspectable.

The project now says what I wanted it to say: I researched imperfect
information and multi-agent learning, built a working game-solving demo,
evaluated it, and published the limitations clearly.

\end{document}
